\section{Техническое задание}

\subsection{Основание для разработки}

Основанием для разработки является задание на выпускную квалификационную работу бакалавра по теме <<Веб-система для автоматизированного управления клиентами интернет-провайдера на базе MikroTik>>. Разработка ведется в рамках преддипломной (производственной) практики, выполняемой на базе университетской кафедры программной инженерии.

Тема и состав работ выпускной квалификационной работы соответствуют направлению подготовки 09.03.04 <<Программная инженерия>>, профилю <<Разработка программно-информационных систем>>. Постановка задачи исходит из проведенного в разделе~1 анализа предметной области, в результате которого установлено отсутствие на рынке готового решения, объединяющего в одной кодовой базе глубокую интеграцию с маршрутизаторами MikroTik по нативному протоколу RouterOS~API, биллинг с автоматическим выпуском счетов и приостановкой обслуживания, подсистему регистрации обращений абонентов и современный веб-интерфейс при разумной совокупной стоимости владения.

Все требования, формулируемые в настоящем разделе, носят функционально-завершенный характер и представляют собой основу для проектирования и реализации программного изделия в разделах~3 и~4.

\subsection{Цель и назначение разработки}

\emph{Целью разработки} является создание веб-ориентированной информационной системы автоматизированного управления абонентами интернет-провайдера малого и среднего масштаба, обеспечивающей сокращение трудозатрат сотрудников на рутинные операции по обслуживанию абонентов и согласованное управление учетной, финансовой и сетевой частями обслуживания.

\emph{Назначение разработки} состоит в предоставлении двух функциональных рабочих мест в рамках единой панели администрирования и общего программного интерфейса прикладного уровня:
\begin{enumerate}
\item Рабочее место сотрудника абонентского отдела~--~для регистрации новых абонентов, ведения договоров, изменения статусов обслуживания, обработки обращений в техническую поддержку, поступивших по внешним каналам связи.
\item Рабочее место сетевого администратора~--~для управления маршрутизаторами MikroTik, правилами брандмауэра и трансляции адресов, мониторинга подключенных устройств, синхронизации тарифных ограничений на стороне маршрутизатора.
\end{enumerate}

В состав разрабатываемой системы входят:
\begin{itemize}
\item единый реестр абонентов с историей подключений, балансом лицевого счета и текущим активным тарифным планом;
\item каталог тарифных планов с параметрами скорости загрузки и отдачи, ежемесячной стоимости, признака симметричности и привязки к именованной очереди маршрутизатора;
\item подсистема выставления счетов с поддержкой расчетного периода, автоматического перевыставления, фиксации фискальных реквизитов на момент выпуска документа и формирования печатной формы счета в формате PDF;
\item подсистема обработки платежей через лицевой счет (кошелек) абонента с поддержкой ручного пополнения сотрудником, оплаты счетов с лицевого счета, фиксации прикрепленного подтверждения платежа в облачном хранилище;
\item подсистема интеграции с маршрутизаторами MikroTik по протоколу RouterOS~API: управление подключенными устройствами, синхронизация простых очередей, чтение и применение правил брандмауэра и трансляции адресов с поддержкой отката изменений, мониторинг беспроводных абонентов;
\item подсистема автоматической приостановки и восстановления обслуживания: фоновая проверка просроченных счетов, перевод абонента в состояние ограничения или приостановки, восстановление обслуживания при погашении задолженности;
\item подсистема технической поддержки в составе панели администрирования для регистрации обращений, поступивших по внешним каналам связи, и ведения внутренней переписки сотрудников по обращению;
\item панель администрирования с ролевой моделью разграничения доступа на основе разрешений в формате \texttt{<<модуль.действие>>};
\item панель аналитики с ключевыми показателями деятельности оператора: численность активных и приостановленных абонентов, объем выставленных и оплаченных средств за период, динамика выручки, перечень крупнейших дебиторов;
\item подсистема системных настроек, позволяющая конфигурировать поведение биллинга, расписание фоновых задач и фискальные реквизиты выпускающей организации без вмешательства в исходный код;
\item подсистема импорта данных из табличных файлов (CSV) с возможностью отката импортированных записей при обнаружении ошибок;
\item подсистема сквозного аудита изменений с фиксацией старого и нового значений изменяемых полей, идентификатора инициатора изменения, IP-адреса источника запроса.
\end{itemize}

\emph{Замечание о границах работы.} Реализация отдельного веб-приложения <<личный кабинет абонента>>, предоставляющего абоненту самостоятельный доступ к данным своего лицевого счета, истории платежей, онлайн-пополнению баланса и самостоятельному созданию обращений в техническую поддержку, в настоящей работе \emph{не выполняется} и рассматривается как направление дальнейшего развития системы (см. заключение). Серверная часть системы и реляционная модель данных при этом проектируются с учетом возможности дальнейшего подключения такого личного кабинета без существенной переработки кодовой базы и схемы базы данных.

\subsection{Требования к функциональности}

Функциональные требования к разрабатываемой системе сгруппированы в двадцать пунктов, охватывающих все рабочие места и подсистемы из подраздела~2.2.

\begin{enumerate}
\item \emph{Авторизация сотрудников.} Вход по электронной почте и паролю (не менее 8 символов с заглавной/строчной буквой и цифрой); выдача токена Bearer~Sanctum для последующих запросов.

\item \emph{Управление учетными записями сотрудников.} CRUD над сотрудниками суперадминистратором с поддержкой мягкого удаления и восстановления до 30 дней; каждый сотрудник имеет ровно одну роль.

\item \emph{Ролевая модель и разрешения.} Разрешения формата \texttt{<<модуль.действие>>}; роли создаются и настраиваются суперадминистратором, роль \texttt{super\_admin} защищена от изменений.

\item \emph{Управление абонентами.} CRUD по карточкам с полями ФИО, документ, контакты, адрес, координаты, IP, дата договора; удаление заменено статусом <<договор расторгнут>>.

\item \emph{Состояние обслуживания.} Четыре статуса (активен, ограничен с HTTP-заглушкой, приостановлен, договор расторгнут); переходы автоматические или ручные.

\item \emph{Управление тарифными планами.} CRUD по тарифам с параметрами скорости, цены, расчетного периода и привязкой к очереди MikroTik; изменения цены не распространяются на существующие подписки.

\item \emph{Подключение к тарифному плану.} Связывание абонента с тарифом с фиксацией текущей цены в подписке и автоматической синхронизацией очереди на маршрутизаторе.

\item \emph{Лицевой счет абонента.} Кошелек с балансом и операциями пополнения, списания, оплаты, возврата с фиксацией в журнале транзакций.

\item \emph{Ручное пополнение лицевого счета.} Зачисление средств сотрудником с прикреплением документа-подтверждения в Cloudinary.

\item \emph{Ручное формирование счета.} Финансовый администратор выпускает счет с автогенерацией фискального номера и фиксацией снимка реквизитов организации.

\item \emph{Автоматическое формирование счетов.} Ежемесячный планировщик создает счета для активных подписок 1-го числа в 00:05 с гарантией идемпотентности.

\item \emph{Автоматическое списание.} После выпуска счета подсистема пытается списать сумму с лицевого счета; при недостатке~--~статус <<ожидает оплаты>>.

\item \emph{Автоматическая приостановка.} Ежедневно в 08:00 абоненты с просрочкой свыше \texttt{grace\_days} (5) переводятся в \emph{ограничен}; свыше \texttt{strict\_days} (15)~--~в \emph{приостановлен}.

\item \emph{Автоматическое восстановление.} При пополнении баланса ограниченного абонента, покрывающем просроченные счета,~--~автоматическая оплата и восстановление обслуживания.

\item \emph{Управление маршрутизаторами MikroTik.} Реестр с зашифрованными паролями (AES-256-CBC, ключ \texttt{APP\_KEY}); чтение пароля через API запрещено.

\item \emph{Синхронизация простых очередей.} Автоматическая (06:00 и 18:00) и ручная синхронизация подписок с \texttt{/queue/simple/} по алгоритму <<вычислить разность~--~применить порциями>>.

\item \emph{Управление брандмауэром и NAT.} Чтение снимка, локальный черновик, валидация, атомарное применение с журналом \texttt{FirewallApplyLog} и откатом на любой предшествующий снимок.

\item \emph{Подсистема технической поддержки.} Регистрация обращений сотрудниками и ведение внутренней переписки по обращению; вложения в Cloudinary; фиксация результата при закрытии.

\item \emph{Системные настройки и автоматизации.} Типизированные настройки в \texttt{system\_settings}, конфигурируемое расписание задач, CSV-импорт с откатом, сквозной аудит изменений (старые/новые значения, инициатор, IP).
\end{enumerate}

\subsection{Требования к интерфейсу}

Интерфейс панели администрирования реализован как одностраничное веб-приложение (SPA) на базе Svelte и метафреймворка SvelteKit. После первоначальной загрузки страница работает без полных перезагрузок: переход между разделами обновляет только рабочую область, шапка и боковая панель остаются на экране без перерисовки.

Поддерживается адаптивная компоновка для устройств с шириной от 360 пикселей (мобильный режим со скрытой боковой панелью) до полноэкранного режима ($\geq 1280$~пикселей). Цветовая схема реализована в темном оформлении на основе темы Cerberus библиотеки Skeleton~UI поверх Tailwind~CSS~v4; все цвета вынесены в CSS-переменные.

Боковая панель навигации содержит пять групп пунктов меню: \emph{Общее} (Главная), \emph{Управление} (Клиенты, Тарифные планы), \emph{Финансы} (Счета, Поставщики), \emph{Инфраструктура} (MikroTik с подменю: Устройства, Брандмауэр, Очереди, Мониторинг, Синхронизация), \emph{Система} (Сотрудники, Конфигурации). При недостаточных правах сотрудника отдельные пункты меню скрываются.

Типовая компоновка реестров сущностей включает заголовок раздела с кнопкой <<Создать>>, строку фильтров и поиска, табличное представление с пагинацией (25 записей на страницу) и боковую панель детального просмотра/редактирования. Для разделов технической поддержки применяется компоновка <<двух колонок>>: список обращений слева, активное обращение с историей переписки справа. Необратимые действия (удаление, расторжение договора, откат брандмауэра) сопровождаются модальным окном подтверждения.

\subsection{Моделирование вариантов использования}

Моделирование вариантов использования выполнено на языке UML согласно спецификации OMG~\cite{omguml}, с учетом методических рекомендаций~\cite{fowler_uml}. В качестве инструмента построения диаграмм использован пакет TikZ для LaTeX, что обеспечивает интеграцию диаграмм непосредственно в исходный код документа без зависимости от внешних графических редакторов.

В разрабатываемой системе предусмотрены четыре действующих лица (актера), все из которых~--~сотрудники интернет-провайдера, входящие в систему через панель администрирования:
\begin{enumerate}
\item \emph{Оператор абонентского отдела}~--~сотрудник, отвечающий за регистрацию новых абонентов, ведение карточек существующих абонентов, обработку обращений в техническую поддержку, поступивших по внешним каналам связи; обладает разрешениями групп \texttt{clientes.*}, \texttt{soporte.*}, \texttt{planes.ver}, \texttt{facturas.ver}.
\item \emph{Финансовый администратор}~--~сотрудник, отвечающий за управление тарифной сеткой, выставление счетов, ручное зачисление средств на лицевые счета абонентов; обладает разрешениями групп \texttt{planes.*}, \texttt{facturas.*}, \texttt{clientes.ver}, \texttt{clientes.editar} в части пополнения лицевого счета.
\item \emph{Сетевой администратор}~--~сотрудник, отвечающий за управление маршрутизаторами MikroTik, правилами брандмауэра, синхронизацией очередей, мониторингом сети; обладает разрешениями групп \texttt{mikrotik.*}, \texttt{configuracion.ver}.
\item \emph{Суперадминистратор}~--~сотрудник, обладающий полным набором разрешений, в том числе правом управления учетными записями сотрудников, ролями и системными настройками; единственный, кто может добавлять и удалять маршрутизаторы в реестре системы и управлять правами доступа.
\end{enumerate}

Иерархия действующих лиц организована по принципу наследования прав: финансовый администратор, оператор абонентского отдела и сетевой администратор наследуют общий набор разрешений на просмотр карточек абонентов и тарифов; суперадминистратор наследует все права остальных ролей. На рисунке~\ref{fig:actors_hierarchy} приведена обобщенная диаграмма действующих лиц. \emph{Замечание.} Абонент (конечный потребитель услуг ИП) в настоящей работе как действующее лицо системы не рассматривается, поскольку реализация личного кабинета абонента вынесена за рамки работы.

\begin{figure}[ht]
\centering
\includegraphics[width=\linewidth,height=0.55\textheight,keepaspectratio]{diagrama_2_1}
\caption{Иерархия действующих лиц системы}
\label{fig:actors_hierarchy}
\end{figure}

Укрупненная диаграмма прецедентов сотрудников с разбиением по ролям приведена на рисунке~\ref{fig:uc_employee}.

\begin{figure}[H]
\centering
\includegraphics[width=\linewidth,height=0.55\textheight,keepaspectratio]{diagrama_2_3}
\caption{Укрупненная диаграмма прецедентов сотрудников}
\label{fig:uc_employee}
\end{figure}

На основании анализа предметной области в системе реализованы следующие тридцать семь прецедентов, группируемые по подсистемам (аутентификация и профили; абоненты; тарифные планы; финансы; маршрутизаторы; брандмауэр и сеть; поддержка и аналитика; администрирование):
\begin{enumerate}
\item Регистрация сотрудника (выполняется суперадминистратором).
\item Вход сотрудника в панель администрирования.
\item Выход из системы.
\item Редактирование собственного профиля.
\item Регистрация нового абонента.
\item Просмотр карточки абонента.
\item Редактирование данных абонента.
\item Приостановка обслуживания абонента вручную.
\item Активация обслуживания абонента вручную.
\item Расторжение договора с абонентом.
\item Создание тарифного плана.
\item Редактирование тарифного плана.
\item Активация и деактивация тарифного плана.
\item Подключение абонента к тарифному плану.
\item Пополнение лицевого счета абонента сотрудником.
\item Просмотр истории платежей абонента.
\item Ручное формирование счета на оплату.
\item Автоматическое формирование счетов по расписанию.
\item Оплата счета с лицевого счета.
\item Ручная отметка оплаты счета.
\item Аннулирование счета.
\item Печать счета в PDF.
\item Просмотр списка маршрутизаторов MikroTik.
\item Добавление маршрутизатора MikroTik.
\item Проверка состояния маршрутизатора MikroTik.
\item Удаление маршрутизатора из реестра.
\item Синхронизация простых очередей.
\item Очистка осиротевших очередей.
\item Получение снимка правил брандмауэра.
\item Применение пакета изменений правил брандмауэра.
\item Откат изменений правил брандмауэра.
\item Импорт правил брандмауэра с маршрутизатора.
\item Просмотр беспроводных абонентов.
\item Регистрация обращения в техническую поддержку (сотрудником со слов абонента, обратившегося по внешнему каналу связи).
\item Добавление внутреннего комментария сотрудника по обращению.
\item Закрытие обращения с фиксацией результата обслуживания.
\item Просмотр панели аналитики.
\item Управление ролями и системными настройками.
\end{enumerate}

Описание прецедентов выполнено по единой схеме, включающей заинтересованные лица и их требования, предусловие, постусловие, основной успешный сценарий, при наличии~--~альтернативные сценарии. Ниже подробно рассмотрены три прецедента, оказывающих наибольшее влияние на проектирование архитектуры системы: создание абонента с автоматической синхронизацией маршрутизатора, изменение статуса обслуживания и автоматическое формирование счетов; типовые операции CRUD сгруппированы в обзорные описания.

\subsubsection{Прецедент <<Регистрация нового абонента>>}

\emph{Заинтересованные лица:} оператор абонентского отдела намерен зарегистрировать в системе нового абонента и подключить его к тарифному плану.

\emph{Предусловие:} оператор аутентифицирован, обладает правом \texttt{clientes.crear}; на маршрутизаторе MikroTik имеется свободный IP-адрес.

\emph{Постусловие:} в системе создана запись абонента, подписки на тарифный план и лицевого счета; на маршрутизаторе создана простая очередь с параметрами тарифа.

\emph{Основной успешный сценарий:}
\begin{enumerate}
\item Оператор открывает раздел <<Клиенты>> и нажимает кнопку <<Создать>>.
\item Система отображает многошаговую форму: личные данные, адрес установки и координаты, технические параметры (IP, выбор тарифа, метод оплаты).
\item Оператор заполняет поля и подтверждает создание.
\item Сервер выполняет валидацию, создает в одной транзакции записи абонента, подписки и кошелька.
\item Сервер диспатчит фоновую задачу синхронизации простых очередей с приоритетом \texttt{high}.
\item Задача подключается к маршрутизатору и создает очередь с именем \texttt{plan-\{plan\_id\}-cli-\{client\_id\}}, целевым IP и максимальной скоростью тарифа.
\item Клиент закрывает форму и переключает контекст на карточку нового абонента.
\end{enumerate}

\emph{Альтернативный сценарий:} при невалидных данных сервер возвращает ответ 422 с детализированным списком ошибок; клиент отображает ошибки рядом с полями формы.

\subsubsection{Прецедент <<Приостановка обслуживания абонента>>}

\emph{Заинтересованные лица:} оператор или подсистема автоматического биллинга намерены приостановить обслуживание абонента.

\emph{Предусловие:} абонент в статусе <<активен>> или <<ограничен>>.

\emph{Постусловие:} абонент переведен в статус <<приостановлен>>; на маршрутизаторе удалена простая очередь абонента и IP добавлен в адресный список \texttt{suspended}; в журнале аудита зафиксирована запись.

\emph{Основной успешный сценарий:} сервер изменяет статус абонента в транзакции, фиксирует событие в таблице \texttt{client\_events}, диспатчит фоновую задачу применения изменений на маршрутизаторе. Задача удаляет очередь и добавляет IP в адресный список приостановленных абонентов.

\subsubsection{Прецедент <<Автоматическое формирование счетов по расписанию>>}

\emph{Заинтересованные лица:} подсистема биллинга действует от имени системного инициатора; цель~--~своевременный выпуск счетов всем абонентам.

\emph{Предусловие:} наступило время срабатывания задачи \texttt{GenerateMonthlyInvoices} согласно расписанию.

\emph{Постусловие:} для каждой активной подписки с наступившей датой следующего выставления создан счет; абонентам с достаточным балансом счета автоматически оплачены.

\emph{Основной успешный сценарий:}
\begin{enumerate}
\item Планировщик Laravel в установленное время вызывает обработчик \texttt{GenerateMonthlyInvoices}.
\item Задача загружает все активные подписки и для каждой проверяет наличие уже выпущенного счета за текущий период (идемпотентность).
\item Для подписок без счета задача формирует новый: генерирует уникальный номер вида \texttt{001-001-NNNNNNNNN}, фиксирует снимок фискальных реквизитов, создает запись \texttt{Invoice}.
\item Задача предпринимает попытку автоматического списания с лицевого счета.
\item Задача формирует сводный отчет о результатах работы в журнале задач.
\end{enumerate}

\subsubsection{Сводное описание остальных прецедентов}

Остальные прецеденты реализуют типовые операции и группируются следующим образом.

\emph{Прецеденты аутентификации и профиля} (1--4): вход сотрудника, выход из системы, редактирование профиля. Реализуются стандартными механизмами Laravel Sanctum: выдача токена при входе, удаление токена при выходе, проверка токена и прав в промежуточных обработчиках для остальных запросов.

\emph{Прецеденты управления абонентами} (6--10): просмотр карточки, редактирование данных, активация, расторжение договора. Реализуются операциями CRUD через эндпоинты \texttt{/api/admin/clientes/...} с обязательной проверкой разрешений формата \texttt{clientes.<действие>} и фиксацией изменений в журнале аудита.

\emph{Прецеденты управления тарифными планами} (11--14): создание, редактирование, активация/деактивация, подключение абонента к тарифу. Особенность: изменение цены тарифа не распространяется на существующие подписки (для них действует цена, зафиксированная в поле \texttt{current\_price} при подключении).

\emph{Прецеденты управления финансами} (16--24): пополнение лицевого счета, просмотр истории платежей, ручное формирование счета, оплата с лицевого счета, ручная отметка оплаты, аннулирование счета, печать в PDF. Печатная форма формируется на стороне клиента средствами jsPDF из снимка фискальных реквизитов; аннулирование возможно только для неоплаченных счетов.

\emph{Прецеденты управления маршрутизаторами и сетью} (25--35): добавление, удаление, проверка состояния маршрутизаторов, синхронизация очередей, очистка осиротевших, получение и откат снимков правил брандмауэра, импорт правил с маршрутизатора, просмотр беспроводных абонентов. Доступны пользователям с разрешениями группы \texttt{mikrotik.*}; пароли маршрутизаторов хранятся в зашифрованном виде ключом \texttt{APP\_KEY}.

\emph{Прецеденты технической поддержки и аналитики} (37--40): регистрация обращения, добавление комментария сотрудника, закрытие обращения с оценкой, просмотр панели аналитики, управление ролями и системными настройками. Аналитические показатели агрегируются с кэшированием в \texttt{cookie} клиента на одни сутки.

\subsection{Требования к надежности}

Целостность данных в системе обеспечивается тремя механизмами.

\emph{Транзакционность операций.} Все операции, затрагивающие несколько таблиц одновременно (создание абонента с одновременным созданием подписки и лицевого счета; создание счета с одновременным списанием средств; применение пакета изменений правил брандмауэра), выполняются в одной транзакции базы данных. Сбой на любом шаге транзакции приводит к откату всех изменений в ее рамках, что исключает появление в базе частично сохраненных или противоречивых записей. Используется механизм \texttt{transaction.atomic} прикладного фреймворка Laravel, оборачивающий блок операций в одну транзакцию с уровнем изоляции \texttt{REPEATABLE READ}.

\emph{Идемпотентность фоновых задач.} Все автоматические задачи спроектированы с учетом возможности повторного запуска: задача автоматического формирования счетов проверяет наличие уже выпущенного счета за расчетный период; задача синхронизации простых очередей применяет операции изменения только при фактическом расхождении состояний; задача автоматической приостановки проверяет текущий статус абонента и не выполняет операцию повторно. Это позволяет безопасно перезапускать задачи в случае сбоев без риска повторного списания средств или дублирующих изменений на маршрутизаторе.

\emph{Восстановление после сбоев интеграции.} Сбой связи с маршрутизатором при выполнении операции синхронизации не приводит к повреждению ни локальной базы данных, ни конфигурации маршрутизатора: задача фиксирует возникшую ошибку в журнале задач, переходит к следующему элементу очереди (если ошибка не критическая) либо приостанавливает выполнение до повторного запуска (если ошибка критическая, например, отказ в аутентификации). Применение изменений правил брандмауэра выполняется с механизмом автоматического отката (см. соответствующий прецедент): при возникновении ошибки на одном из шагов система автоматически возвращает маршрутизатор к состоянию <<до>>, зафиксированному в журнале применений.

Резервное копирование базы данных осуществляется штатными средствами СУБД MySQL (\texttt{mysqldump}) либо средствами специализированных пакетов прикладного фреймворка (\texttt{spatie/laravel-backup}), с настраиваемой периодичностью и сроком хранения копий. Восстановление выполняется из любой имеющейся копии без дополнительного инструментария.

При плановой остановке сервера для обновления программного обеспечения система отображает на странице входа панели администрирования и личного кабинета сообщение о технических работах с указанием планового времени восстановления.

\subsection{Требования к информационной безопасности}

Аутентификация всех запросов, кроме точек входа и получения CSRF-токена, выполняется через токены Bearer библиотеки Laravel~Sanctum~\cite{sanctum}, хранящиеся в таблице \texttt{personal\_access\_tokens}. Пароли учетных записей хранятся в виде хешей Bcrypt с 12 раундами. Пароли маршрутизаторов MikroTik в реестре системы шифруются симметрично AES-256-CBC с ключом \texttt{APP\_KEY}; чтение в открытом виде через программный интерфейс не предусмотрено.

Доступ разграничивается ролевой моделью с разрешениями формата \texttt{<<модуль.действие>>}; проверка выполняется в промежуточных обработчиках Laravel до передачи запроса в контроллер. Недостаточный уровень прав приводит к ответу 403 (Forbidden). Доступ к данным конкретного абонента имеют только сотрудники с соответствующими правами и сам абонент.

Защита от внедрения SQL-кода обеспечивается единым использованием Eloquent~ORM с параметризованными запросами~\cite{owasp}; конкатенация пользовательского ввода в SQL-выражения в коде запрещена. Защита от XSS выполняется автоматическим экранированием в шаблонизаторе Svelte, защита от CSRF~--~через cookie csrftoken и заголовок \texttt{X-CSRFToken}. Все изменения данных автоматически фиксируются в журнале аудита через trait \texttt{Auditable} с указанием инициатора, IP-адреса источника и старого/нового значений измененных полей.

\subsection{Требования к программной и технической совместимости}

\emph{Серверная часть}: Linux (Ubuntu~22.04~LTS+ или Debian~12+) либо Windows Server~2019+; PHP~8.2+ с расширениями \texttt{pdo\_mysql}, \texttt{mbstring}, \texttt{openssl}, \texttt{bcmath}, \texttt{gd}, \texttt{intl}, \texttt{sockets}; Composer~2+; Node.js~20.19+ (сборка клиента); MySQL~8.0+ или MariaDB~10.6+; Nginx~1.18+ или Apache~HTTPD~2.4+; опционально Redis~7+ для очередей задач.

\emph{Аппаратные требования к серверу} (для базы до 5000 абонентов): процессор 2~ядра $\geq$~2~ГГц, ОЗУ 4~ГБ, диск 40~ГБ, постоянное подключение к сети, статический внешний IP-адрес.

\emph{Клиентская часть}: любой современный браузер с поддержкой ECMAScript~2020: Chrome~120+, Firefox~120+, Edge~120+, Safari~17+. Минимальное разрешение~--~360~пикселей по ширине. Поддержка Internet~Explorer не предусмотрена.

\emph{Сетевые требования}: TCP-порт 8728 (или 8729 для TLS) между сервером и маршрутизаторами MikroTik; TCP-порт 443 (HTTPS) между клиентами и сервером. В производственной конфигурации канал шифруется TLS~1.3 с сертификатом Let's~Encrypt; HTTP перенаправляется на HTTPS.

\emph{Стек серверной части}: PHP~8.2+, Laravel~12~\cite{laraveldoc}, Laravel~Sanctum~4~\cite{sanctum}, \texttt{evilfreelancer/routeros-api-php}~1.6~\cite{routerosphp}, \texttt{cloudinary-labs/cloudinary-laravel}~3.0, PestPHP~4.1, MySQL~8.0~\cite{mysqldoc}.

\emph{Стек клиентской части}: Svelte~5.39~\cite{sveltedoc}, SvelteKit~2.43~\cite{sveltekitdoc}, Tailwind~CSS~4.1~\cite{tailwind}, Skeleton~UI~4.7, jsPDF~3.0 + jspdf-autotable~5.0, Lucide~Icons~0.546, svelte-sonner~1.0, recharts~3.3, Vite~7.1, Vitest~4.1.

\subsection{Требования к оформлению документации}

Разработка программной документации и программного изделия должна производиться согласно ГОСТ~19.102--77~\cite{gost19102} (стадии разработки программ и программной документации) и ГОСТ~34.601--90~\cite{gost34601} (комплекс стандартов на автоматизированные системы, стадии создания). Качество разрабатываемого программного обеспечения оценивается по модели характеристик качества, установленной ГОСТ~Р~ИСО/МЭК~25010--2015~\cite{gost25010} (модели качества систем и программных продуктов).

Оформление настоящей пояснительной записки выполнено в соответствии с СТУ~02.030--2023~\cite{stu02030} (стандарт университета на оформление курсовых работ и выпускных квалификационных работ) и ГОСТ~7.32--2017~\cite{gost7320} (структура и правила оформления отчета о научно-исследовательской работе).

Список использованных источников оформлен в соответствии с ГОСТ~Р~7.0.100--2018~\cite{gost70100} (правила составления библиографического описания). Ссылки в тексте на источники выполнены в квадратных скобках с указанием порядкового номера источника по списку использованных источников.
